% Part: normal-modal-logic
% Chapter: frame-definability
% Section: definability

\documentclass[../../../include/open-logic-section]{subfiles}

\begin{document}

\olfileid{nml}{frd}{def}

\olsection{ફ્રેમ-વ્યાખ્યાયિતતા}

\olref[acc]{thm:soundschemas}નાં વિપરીત વિધાનો
નિદર્શનો માટે સાચાં નથી, છતાં ``નિદર્શન''ને બદલે
``ફ્રેમ'' મૂકીએ તો તે સાચાં છે: \olref[acc]{thm:soundschemas}માં
વિચારેલા ગુણધર્મો માટે !!a{formula} કોઈ \emph{ફ્રેમ}માં
માન્ય હોય તો તે ફ્રેમનો સુલભતા સંબંધ અનુરૂપ
ગુણધર્મ ધરાવે છે. તેથી વિચારેલાં !!{formula}s તે
ગુણધર્મ ધરાવતી ફ્રેમોના વર્ગોને \emph{વ્યાખ્યાયિત}
કરે છે.

\begin{defn}
  $\mClass{F}$ ફ્રેમોનો વર્ગ હોય, તો $!A$ એ
  $\mClass{F}$ને \emph{વ્યાખ્યાયિત કરે છે} એમ કહીએ
  ત્યારે અને માત્ર ત્યારે જ જ્યારે બધી અને માત્ર
  બધી ફ્રેમો~$\mModel{F} \in \mClass{F}$ માટે
  $\mModel{F} \Entails !A$ હોય.
\end{defn}

હવે ફ્રેમો માટે વ્યાખ્યાયિતતાના પૂર્ણ પરિણામો
સ્થાપીએ.

\begin{thm}\ollabel{thm:fullCorrespondence}
\olref[acc]{tab:five}ની જમણી બાજુનું !!{formula}
ફ્રેમ~$\mModel{F}$માં માન્ય હોય, તો $\mModel{F}$
ડાબી બાજુનો ગુણધર્મ ધરાવે છે.
\end{thm}

\begin{proof}
  \begin{enumerate}
  \item માનો કે \Ax{D} એ $\mModel{F} = \tuple{W, R}$માં
    માન્ય છે, એટલે કે
    $\mModel{F} \Entails \Box p \lif \Diamond p$.
    ફ્રેમ~$\mModel{F}$ પર આધારિત નિદર્શન $\mModel{M} =
    \tuple{W, R, V}$ અને $w \in
    W$ લો. એવું કોઈ $v$ છે કે~$Rwv$ તે બતાવવું છે.
    નથી એમ માનો: તો કોઈપણ~$!A$ માટે
    $\mSat{M}{\Box !A}[w]$ અને
    $\mSat/{M}{\Diamond !A}[w]$ બંને છે, ખાસ કરીને~$p$
    માટે. પરંતુ ત્યારે $\mSat/{M}{\Box p \lif
      \Diamond p}[w]$, જે $\mModel{F}
    \Entails \Box p \lif \Diamond p$ની ધારણાથી વિરોધાભાસ છે.
  \item માનો કે \Ax{T} એ $\mModel{F}$માં માન્ય છે,
    એટલે કે $\mModel{F} \Entails \Box
    p \lif p$. $w \in W$ કોઈપણ જગત લો;
    $Rww$ બતાવવું છે. $u \in V(p)$ ત્યારે અને માત્ર
    ત્યારે જ મૂકો જ્યારે $Rwu$ હોય ($q$ એ $p$ સિવાય
    બીજું ચલ હોય ત્યારે $V(q)$ ગમે તેમ લો, જેમ કે
    $V(q) = \emptyset$). હવે
    $\mModel{M} = \tuple{W, R, V}$ લો. રચના મુજબ
    $Rwu$ ધરાવતા બધા $u$ માટે $\mSat{M}{p}[u]$,
    તેથી $\mSat{M}{\Box p}[w]$. પરંતુ ધારણા મુજબ
    $\Box p \lif p$ એ~$w$માં સત્ય છે, તેથી
    $\mSat{M}{p}[w]$; અને $V$ની વ્યાખ્યા મુજબ આ
    માત્ર~$Rww$ હોય તો જ શક્ય છે.
  \item પ્રતિધનાત્મક રીતે સાબિત કરીએ: માનો કે
    $\mModel{F}$ સંમિત નથી; તો \Ax{B}, એટલે કે
    $p \lif \Box\Diamond p$, એ $\mModel{F}= \tuple{W, R}$માં
    માન્ય નથી તે બતાવીએ. $\mModel{F}$ સંમિત ન હોય,
    તો એવાં $u$, $v \in W$ છે કે $Ruv$ છે પણ
    $Rvu$ નથી. $V$ એવું વ્યાખ્યાયિત કરો કે
    $w \in V(p)$ ત્યારે અને માત્ર ત્યારે જ જ્યારે
    $Rvw$ ન હોય (બાકીનાં માટે $V$ ગમે તેમ હોય).
    $\mModel{M} =\tuple{W, R,
      V}$ લો. હવે $V$ની વ્યાખ્યા મુજબ $Rvw$ ન હોય
    એવા બધા $w$ માટે $\mSat{M}{p}[w]$; ખાસ કરીને,
    $Rvu$ ન હોવાથી $\mSat{M}{p}[u]$. વળી,
    $Rvw$ ત્યારે અને માત્ર ત્યારે જ જ્યારે
    $w \notin V(p)$, તેથી $Rvw$ અને
    $\mSat{M}{p}[w]$ બંને હોય એવું કોઈ~$w$ નથી;
    પરિણામે $\mSat/{M}{\Diamond p}[v]$. $Ruv$
    પણ હોવાથી $\mSat/{M}{\Box\Diamond p}[u]$.
    તેથી $\mSat/{M}{p \lif
      \Box\Diamond p}[u]$, અને આમ $\Ax{B}$
    ફ્રેમ~$\mModel{F}$માં માન્ય નથી.
  \item માનો કે \Ax{4} એ $\mModel{F} = \tuple{W,R}$માં
    માન્ય છે, એટલે કે
    $\mModel{F} \Entails \Box p \lif \Box\Box p$;
    અને $u$, $v$, $w \in W$ એવાં કોઈપણ જગતો લો
    કે $Ruv$ અને $Rvw$. $Ruw$ બતાવવું છે.
    $V$ એવું વ્યાખ્યાયિત કરો કે $z \in V(p)$ ત્યારે
    અને માત્ર ત્યારે જ જ્યારે $Ruz$ હોય (બાકીનાં
    માટે $V$ ગમે તેમ હોય). હવે
    $\mModel{M} =\tuple{W, R, V}$ લો. $V$ની વ્યાખ્યા
    મુજબ $Ruz$ ધરાવતા બધા $z$ માટે
    $\mSat{M}{p}[z]$, તેથી $\mSat{M}{\Box p}[u]$.
    પરંતુ ધારણા મુજબ \Ax{4}, એટલે કે $\Box p
    \lif \Box \Box p$, એ $u$માં સત્ય છે, તેથી
    $\mSat{M}{\Box \Box
      p}[u]$. $Ruv$ અને $Rvw$ હોવાથી
    $\mSat{M}{p}[w]$; $V$ની વ્યાખ્યા પ્રમાણે આ
    માત્ર $Ruw$ હોય તો જ શક્ય છે, જેમ જોઈએ તેમ.
  \item પ્રતિધનાત્મક રીતે આગળ વધીએ: ફ્રેમ
    $\mModel{F} = \tuple{W, R}$ યુક્લિડિયન નથી એમ
    માની તે \Ax{5}ને ખોટું પાડે છે, એટલે કે
    $\mModel{F} \Entails/ \Diamond p \lif
      \Box\Diamond p$, તે બતાવીએ. એવાં જગતો
    $u$, $v$, $w \in W$ માનો કે $Rwu$ અને $Rwv$
    છે, પણ $Ruv$ નથી. $V$ એવું વ્યાખ્યાયિત કરો
    કે બધાં જગતો $z$ માટે $z \in V(p)$ ત્યારે અને
    માત્ર ત્યારે જ જ્યારે~$Ruz$ \emph{ન હોય}.
    $\mModel{M} =\tuple{W, R, V}$ લો. ત્યારે ધારણા
    મુજબ $\mSat{M}{p}[v]$ અને $Rwv$ હોવાથી
    $\mSat{M}{\Diamond p}[w]$ પણ છે. જોકે
    $Ruy$ અને $\mSat{M}{p}[y]$ બંને હોય એવું કોઈ
    જગત $y$ નથી, તેથી
    $\mSat/{M}{\Diamond
      p}[u]$. $Rwu$ હોવાથી પછી
    $\mSat/{M}{\Box\Diamond
      p}[w]$, એટલે \Ax{5}, $\Diamond p \lif \Box\Diamond p$,
    એ~$w$માં ખોટું પડે છે.
  \end{enumerate}
\end{proof}
\sourcecorrection{OLMD-015}{સ્થિર અંગ્રેજી મૂળમાં
\Ax{D}ના પ્રમાણની એક સંતોષ-અભિવ્યક્તિમાં
જગતનો નિર્દેશ છૂટી ગયો છે; અન્ય જગત-સાપેક્ષ
અભિવ્યક્તિઓની જેમ ત્યાં પણ નિર્દેશ પૂર્યો છે.}
\sourcecorrection{OLMD-016}{\Ax{T}ના પ્રમાણમાં
સ્થિર અંગ્રેજી મૂળે $V(q)$ના ઉદાહરણનું બંધ
કૌંસ ગણિત-અભિવ્યક્તિની અંદર મૂકે છે;
તેને ગણિત-અભિવ્યક્તિની બહાર મૂક્યું છે.}

\Ax{D}ના પ્રમાણ અને અન્ય કિસ્સાઓ વચ્ચેનો એક
ભેદ ધ્યાનમાં આવશે: નિમણૂક~$V$નો તેમાં કોઈ
ઉલ્લેખ થયો નથી. હકીકતમાં, આપણે સાબિત કર્યું
કે $\mSat{M}{\Ax{D}}$ હોય તો $\mModel{M}$ સિરિયલ
છે. તેથી \Ax{D} માત્ર ફ્રેમોનો નહીં, પરંતુ સિરિયલ
\emph{નિદર્શનોનો} વર્ગ વ્યાખ્યાયિત કરે છે.

\begin{cor}\ollabel{prop:D-serial}
  જે નિદર્શનમાં \Ax{D} સત્ય હોય તે સિરિયલ છે.
\end{cor}

\begin{cor}
\olref[acc]{tab:five}ની જમણી બાજુનું દરેક
!!{formula} ડાબી બાજુનો ગુણધર્મ ધરાવતી
ફ્રેમોનો વર્ગ વ્યાખ્યાયિત કરે છે.
\end{cor}

\begin{proof}
  \olref[acc]{thm:soundschemas}માં આપણે સાબિત કર્યું
  કે નિદર્શન ડાબી બાજુનો ગુણધર્મ ધરાવે તો જમણી
  બાજુનું !!{formula} તેમાં સત્ય છે. તેથી ફ્રેમ~$\mModel{F}$
ડાબી બાજુનો ગુણધર્મ ધરાવે તો જમણી બાજુનું
  !!{formula} એ~$\mModel{F}$માં માન્ય છે.
  \olref{thm:fullCorrespondence}માં આપણે વિપરીત
  વિધાનો સાબિત કર્યાં: જમણી બાજુનું !!a{formula}
  એ~$\mModel{F}$માં માન્ય હોય તો $\mModel{F}$ ડાબી
  બાજુનો ગુણધર્મ ધરાવે છે.
\end{proof}

\begin{prob}
જો \olref[nml][frd][acc]{tab:anotherfive}ની જમણી
બાજુનું !!{formula} ફ્રેમ~$\mModel{F}$માં માન્ય હોય,
તો $\mModel{F}$ ડાબી બાજુનો ગુણધર્મ ધરાવે છે તે
બતાવો. તે માટે ડાબી બાજુનો ગુણધર્મ \emph{ન} ધરાવતી
ફ્રેમ લો અને યોગ્ય~$V$ વ્યાખ્યાયિત કરો જેથી
જમણી બાજુનું !!{formula} કોઈ જગતમાં અસત્ય થાય.
\end{prob}

\olref{thm:fullCorrespondence} એ પણ બતાવે છે કે
ગુણધર્મો ભેગા કરી શકાય: દાખલા તરીકે, \Ax{B}
અને \Ax{4} બંને~$\mModel{F}$માં માન્ય હોય તો ફ્રેમ
સંમિત અને પરંપરિત બંને છે, વગેરે. કેટલાંક
ફ્રેમ-ગુણધર્મો ભેગાં ધરાવતી બધી ફ્રેમોમાં માન્ય
!!{formula}sના ગણ તરીકે ઘણાં મહત્ત્વનાં મોડલ
તર્કશાસ્ત્રો લક્ષણિત થાય છે; તેથી અનુરૂપ વ્યાખ્યાયિત
કરતાં !!{formula}s માન્ય હોય તેવી બધી ફ્રેમોમાં
માન્ય !!{formula}sના ગણ તરીકે પણ તેમને લક્ષણિત
કરી શકીએ. દાખલા તરીકે, શાસ્ત્રીય પ્રણાલી
\Log{S4} એ બધી સ્વવાચક અને પરંપરિત ફ્રેમોમાં
માન્ય !!{formula}sનો ગણ છે, એટલે કે \Ax{T}
અને~\Ax{4} બંને માન્ય હોય એવી બધી ફ્રેમોમાં
માન્ય સૂત્રોનો ગણ. \Log{S5} એ બધી સ્વવાચક,
સંમિત અને યુક્લિડિયન ફ્રેમોમાં માન્ય સૂત્રોનો ગણ છે,
એટલે કે \Ax{T}, \Ax{B} અને \Ax{5} ત્રણેય
માન્ય હોય એવી બધી ફ્રેમોમાં માન્ય સૂત્રોનો ગણ.

$R$ના ગુણધર્મો વચ્ચેના તાર્કિક સંબંધો સામાન્ય રીતે
તેમને વ્યાખ્યાયિત કરતાં !!{formula}s વચ્ચેના સંબંધોને
અનુરૂપ હોય છે. દાખલા તરીકે, દરેક સ્વવાચક સંબંધ
સિરિયલ છે; તેથી ફ્રેમમાં \Ax{T} માન્ય હોય તો~\Ax{D}
પણ માન્ય છે. (ધ્યાન આપો કે આ સંબંધ
\emph{ફલિતતાનો નથી}. $\mSat{M}{\Ax{T}}[w]$ હોય
ત્યારે હંમેશાં $\mSat{M}{\Ax{D}}[w]$ હોય એવું નથી.)
આવા કેટલાક સંબંધો નોંધીએ.

\begin{prop}\ollabel{prop:relation-facts}
  $R$ એ ગણ $W$ પરનો દ્વિસ્થાની સંબંધ હોય; તો:
  \begin{enumerate}
  \item $R$ સ્વવાચક હોય તો તે સિરિયલ છે.
  \item $R$ સંમિત હોય, તો તે પરંપરિત ત્યારે અને માત્ર
    ત્યારે જ છે જ્યારે તે યુક્લિડિયન હોય.
  \item $R$ સંમિત અથવા યુક્લિડિયન હોય તો તે
    નબળે દિશિત છે (તેમાં ``હીરા ગુણધર્મ'' છે).
  \item $R$ યુક્લિડિયન હોય તો તે નબળે જોડાયેલો છે.
  \item $R$ વિધેયાત્મક હોય તો તે સિરિયલ છે.
  \end{enumerate}
\end{prop}

\begin{prob}
  \olref[nml][frd][def]{prop:relation-facts} સાબિત કરો.
\end{prob}

\end{document}
